\section{Introduction}\label{sec:intro}

Parton distribution functions (PDFs) are crucial quantities characterizing the structure of hadrons at high energy. They are defined as momentum distributions of quarks and gluons in an infinite-momentum hadron, and provide a necessary input for the description of experimental data from hadron-hadron or lepton-hadron colliders. However, calculating PDFs from first principles has been a long-standing problem in hadron physics, since PDFs are low-energy properties of a hadron defined in terms of lightcone correlations. Traditional lattice QCD methods focus on the calculation of their moments, but can only determine the first few moments due to the power divergent mixing between different moment operators~\cite{Martinelli:1987zd,Martinelli:1988xs,Detmold:2001dv,Dolgov:2002zm}. Reconstructing PDFs, however, requires information on all their moments.

In the past few years, a new approach has been proposed to circumvent the above difficulty, which has now been formulated as the large momentum effective theory (LaMET)~\cite{Ji:2013dva,Ji:2014gla}. LaMET provides a possibility to study parton properties of the hadron in general. For a given parton observable such as PDFs, one can deconstruct the lightcone correlations such that the resulting quantity, known as a quasi-observable, becomes time-independent though no longer frame-independent. The quasi-observable approaches the original parton observable under an infinite Lorentz boost, and has the same infrared (IR) behavior as the latter by construction. In a hadron with finite but large momentum, the quasi-observable can be factorized into the parton observable convoluted with a hard coefficient, up to power corrections suppressed by the hadron momentum~\cite{Ji:2013dva,Ji:2014gla,Xiong:2013bka,Ma:2017pxb,Izubuchi:2018srq,Ji:2015qla,Xiong:2015nua}.
Since LaMET was proposed, much progress has been achieved with respect to both the theoretical understanding of the formalism and the direct calculation of PDFs from lattice QCD (see a recent review~\cite{Cichy:2018mum} and references therein). The lattice calculations have been predominantly focused on isovector quark PDFs, since they do not mix with gluon PDFs. Despite limited volumes and relatively coarse lattice spacings, the state-of-the-art nucleon isovector quark PDFs determined from lattice data at the physical point already show a reasonable agreement~\cite{Chen:2018xof,Lin:2018qky,Alexandrou:2018pbm,Alexandrou:2018eet,Liu:2018hxv} with phenomenological results extracted from experimental data~\cite{Dulat:2015mca,Ball:2017nwa,Harland-Lang:2014zoa,Nocera:2014gqa,Ethier:2017zbq}.

Lattice calculations of gluon parton physics are even more crucial because gluons play an extremely important role in experiments performed at the Large Hadron Collider (LHC) and the future Electron-Ion Collider (EIC) in US. However,
in contrast to the intensive work on extracting quark PDFs using LaMET, much less effort has been devoted to the study of gluon quasi-PDFs.
In Refs.~\cite{Wang:2017qyg,Wang:2017eel}, the first effort was made in understanding the renormalization property of gluon quasi-PDFs, where it was discovered that the gluon quasi-PDF suffers from a mixing with other gluon operators allowed by symmetries of the theory. While on the other hand, in the first exploratory attempt to calculate the gluon unpolarized PDF~\cite{Fan:2018dxu}, the authors assumed a multiplicative renormalization of the gluon quasi-PDF used in their lattice calculation.

In this Letter, we perform a systematic study of the renormalization of gauge-invariant non-local operators defining the gluon quasi-PDFs, focusing on their
power divergence structure. Following an earlier work involving two of us~\cite{Ji:2017oey} (see also~\cite{Ishikawa:2017faj,Green:2017xeu}), we introduce a general framework in which the Wilson line in the adjoint representation is replaced by a product of two auxiliary ``heavy quark" fields. Such a theory with an auxiliary ``heavy quark" can be shown to be renormalizable to all orders in perturbation theory.
We then show how to construct a multiplicatively renormalizable gluon quasi-PDF operator, and explain how to calculate the renormalization factors non-perturbatively on lattice. The analysis is also extended to the polarized case. At last, we present a factorization formula for the renormalized gluon quasi-PDF, from which one can extract the gluon PDF. Our findings remove the obstacle to define a continuous gluon quasi-PDF through lattice simulations.
